\chapter{Каноничность ориентированной конструкции высоты--Ходжа в KO6}

Предыдущая проверка получила точное тождество после введения целочисленной
высоты. Решающий вопрос --- выводится ли эта высота из полной 18-мерной
вещественной поколенческой геометрии или является дополнительным выбором.

\section{Исправления исходной постановки}

В проекте используются знаки KO6
\begin{equation}
 J^2=+1,\qquad JD=DJ,qquad J\Gamma=-\Gamma J,
\end{equation}
в согласии со стандартной таблицей знаков; см. \path{arXiv:1409.5983} и
\path{arXiv:1601.00219}.

Центр прямой суммы способен различать её слагаемые. Рассматриваемая высота
является блочно-скалярной на вершинах неприводимых бимодулей и коммутирует с
левым и противоположным действиями алгебры. Она не обязана принадлежать
образу одного только левого центра.

Предложенное условие
\begin{equation}
 [[h,\pi(a)],J\pi(b)J^{-1}]=0
 \label{eq:v5-height-proposed-first-order}
\end{equation}
не выбирает высоту: для всякого блочно-скалярного \(h\) имеем
\([h,\pi(a)]=0\). Условие первого порядка ограничивает \(D\), а не
дополнительную коммутирующую координату вершины.

\section{Полный набор высот KO6}

Пусть на частичной цепочке
\begin{equation}h_p=\operatorname{diag}(aI_3,bI_3,cI_3).\end{equation}
Естественное нечётное по \(J\) дополнение имеет вид
\begin{equation}
 h_F=h_p\oplus(-h_p).
 \label{eq:v5-full-ko6-height}
\end{equation}
Для любых вещественных \(a,b,c\) оно самосопряжённо, коммутирует с
\(\Gamma_F\) и обоими действиями алгебры и удовлетворяет
\(Jh_FJ^{-1}=-h_F\). Таким образом, условия KO6 оставляют три непрерывных
параметра до нормировки рёбер.

Требование единичных модулей перепадов высоты даёт
\begin{equation}
 (b-a)^2=1,\qquad(c-b)^2=1,
 \label{eq:v5-unit-height-gaps}
\end{equation}
или \(\operatorname{ad}_{h_F}^2(D_F)=D_F\). Даже если дополнительно
потребовать, чтобы \(P_G=I-h_F^2\) был проектором на обе средние копии,
остаются два неэквивалентных класса с точностью до общего обращения знака.

\section{Явное препятствие двух высот}

На частичной цепочке выберем
\begin{align}
 h_{\text{цеп}}&=\operatorname{diag}(-I_3,0_3,+I_3),\\
 h_{\text{сток}}&=\operatorname{diag}(-I_3,0_3,-I_3),
\end{align}
и дополним их по формуле \eqref{eq:v5-full-ko6-height}. Обе высоты проходят
все перечисленные условия и имеют один средний проектор
\begin{equation}
 I-h_{\text{цеп}}^2=I-h_{\text{сток}}^2
 =\operatorname{diag}(0_3,I_3,0_3)
\end{equation}
на частичной половине, но не связаны общим знаком.

Для
\begin{equation}d_h=\frac12(D_F+[h_F,D_F])\end{equation}
согласованная цепочка даёт
\begin{equation}
 [d_{\text{цеп}},d_{\text{цеп}}^\dagger]\big|_{V_G}
 =XX^\dagger-Y^\dagger Y,
 \label{eq:v5-chain-height-difference}
\end{equation}
тогда как высота со стоком в средней вершине даёт
\begin{equation}
 [d_{\text{сток}},d_{\text{сток}}^\dagger]\big|_{V_G}
 =XX^\dagger+Y^\dagger Y.
 \label{eq:v5-sink-height-sum}
\end{equation}

\begin{theorem}[Неединственность высоты]
Вещественность, градуировка, совместимость с обоими действиями алгебры,
предложенное условие первого порядка, единичные перепады и требование к
среднему проектору не определяют согласованную высоту колчана. Существуют по
меньшей мере два неэквивалентных класса, дающих разные знаки в отображении
момента.
\end{theorem}

\section{Нормировка следа}

Удвоение KO6 не вводит независимого относительного веса. Для согласованной
высоты
\begin{equation}
 \Tr_{\mathcal H_F}\left(P_G[d_h,d_h^\dagger]^2\right)
 =2\Tr_3(XX^\dagger-Y^\dagger Y)^2.
\end{equation}
После деления на полную размерность среднего сектора, равную шести,
\begin{equation}
 \frac16\Tr_{\mathcal H_F}\left(P_G[d_h,d_h^\dagger]^2\right)
 =\tau_3(XX^\dagger-Y^\dagger Y)^2.
\end{equation}
Следовательно, нормировка правильна, но ориентация не выбирается.

\section{Связь с литературой}

Разность входящих и исходящих вкладов является стандартным отображением
момента уже ориентированного колчана; см. \path{arXiv:0807.4734}. Функции
высоты, задающие ориентацию колчанов Дынкина, также выступают фиксированными
комбинаторными данными; см. \path{arXiv:2410.10070}. Диаграммы Краевского
несут ориентированные данные представления, но обычная конечная спектральная
тройка сама по себе не порождает конкретную целую высоту; ср.
\path{arXiv:hep-th/0501181}.

\begin{nogocriterion}[Скрытый селектор ориентации]
Тождество \(h^3=h\) не помогает: обе высоты ему удовлетворяют. Требование
трёх различных уровней или монотонности вдоль заранее названной цепочки
ввело бы желаемую ориентацию как новую аксиому.
\end{nogocriterion}

\section{Вывод}

\begin{protocol}
\begin{enumerate}
 \item алгебраическое тождество высоты--Ходжа: \textbf{выполнено};
 \item совместимость согласованной высоты с KO6: \textbf{выполнена};
 \item нормировка полного следа: \textbf{выполнена};
 \item единственность высоты: \textbf{не выполнена};
 \item вывод ориентации из текущей геометрии KO6:
 \textbf{направление закрыто};
 \item физическое замыкание: \textbf{не достигнуто}.
\end{enumerate}
\end{protocol}

Следующая независимая проверка ---
\path{version5_twisted_family_automorphism_gate}.
Машинный сертификат сохранён в
\path{s2t_v5_oriented_height_hodge_ko6_gate_results.json}.